\documentclass[11pt]{article}
\usepackage[margin=1in]{geometry}
\usepackage{amsmath,amssymb,amsthm}
\usepackage[colorlinks=true,linkcolor=blue,urlcolor=blue,citecolor=blue]{hyperref}
\newtheorem{theorem}{Theorem}
\newtheorem{lemma}[theorem]{Lemma}
\newtheorem{corollary}[theorem]{Corollary}
\newtheorem{remark}[theorem]{Remark}
\newcommand{\F}{\mathbb F}
\newcommand{\rank}{\operatorname{rank}}
\newcommand{\nnz}{\operatorname{nnz}}
\newcommand{\one}{\boldsymbol1}
\newcommand{\E}{\mathbb E}
\newcommand{\Prob}{\mathbb P}
\title{Certified full finite-field kernels and expected nearly quadratic partition listing}
\author{Partitioning dynamic session: full-kernel algorithm audit}
\date{November 15, 2026}
\begin{document}
\maketitle
\begin{abstract}
For a sparse symmetric matrix over a fixed prime field, we give an
always-correct Las Vegas algorithm computing and certifying its full
base-field kernel in expected $\widetilde O(n^2(r+1))$ bit work,
where $r$ is its corank. No corank-one promise is made.
The ingredients are classical rank-deficient diagonal preconditioning,
scalar Wiedemann recurrences, and finite-field arithmetic. A verified
kernel basis and a shifted-recurrence rank lower bound certify each
other. An exponential corank moment then yields expected nearly
quadratic exact listing of all balanced sign-kernel partitions,
including instances with extra real eigenspace dimensions. The graph
moment is an explicit input to the application, rather than an
unverified consequence of qualitative contiguity.
\end{abstract}

\section{Unconditional per-matrix algorithm}
Fix an odd prime $p$. Let $M=M^T\in\F_p^{n\times n}$, $n\ge1$,
be supplied as a sparse matrix, with $w=\nnz(M)$,
$k=\rank_{\F_p}M$, and $r=n-k$.
All logarithmic constants below may depend on the fixed $p$.

\begin{theorem}\label{thm:kernel}
There is a Las Vegas algorithm, valid for every such fixed $M$,
which returns a basis $C\in\F_p^{n\times r}$ of its entire kernel,
pivot rows $J$ with $C_{J,:}=I_r$, and a deterministically checkable
rank certificate. Its expected work is
\[
 \widetilde O_p\bigl(n(w+n)(r+1)\bigr)
 \quad\text{bit operations}.
\]
Its workspace is $\widetilde O_p(n(r+1)+w)$ bits, apart from a retry
counter. The certificate has $O_p(nr+n\log n)$ bits and can be
verified in $\widetilde O_p(n(w+n))$ bit work.
For bounded-degree graph matrices, $w=O(n)$, giving expected
$\widetilde O_p(n^2(r+1))$ bit work. These assertions include the
zero matrix and nonsingular matrices. No numerical gap or generic
rank profile is assumed.
\end{theorem}

\section{Extension field and primary preconditioning inputs}
Choose the least positive $s$ for which
\[
 q=p^s\ge64(n+1)^6,
\]
and construct $K=\F_q=\F_p[z]/(h(z))$ once, reusing it across trials,
using a verified monic irreducible polynomial of degree $s$.
Rank and corank are unchanged
by this scalar extension.
In a fresh trial draw independent
\[
 D=\operatorname{diag}(d_1,\ldots,d_n),\quad d_i\text{ uniform in }K^*,
 \quad u,v\text{ uniform in }K^n,\qquad T=DM.
\]
The following are distinct primary results, both valid in the fixed
characteristic $p$, however small relative to $n$.
Chen et al.\ \cite[Theorem 4.2]{chen} give
\begin{equation}\label{eq:cyc}
 \Prob(T\text{ not cyclic on its nonzero primary part})
 \le\frac{n(n-1)}{2(q-1)}.
\end{equation}
Their Theorem 4.5, using symmetry of $M$, gives
\begin{equation}\label{eq:nil}
 \Prob(T\text{ has a nilpotent zero block of size }>1)
 \le\frac r{q-1}.
\end{equation}
No independence between these two events is required for their union
bound. The squarefree conclusion of their Theorem 4.3 has a
characteristic restriction and is \emph{not} invoked.

On their common good event,
\[
 K^n=\ker T\oplus\operatorname{im}T,
 \qquad T|_{\operatorname{im}T}\text{ invertible and cyclic of dimension }k.
\]
Repeated nonzero roots of its minimal polynomial cause no problem.

\section{A certifying trial, including unlucky choices}
Compute exactly the $2n$ scalars
\begin{equation}\label{eq:seq}
 a_j=u^TT^{j+1}v,\qquad 0\le j<2n.
\end{equation}
Use Berlekamp--Massey to obtain their true shortest recurrence length
$\ell$ and monic polynomial $g$, of degree $\ell$.
If the implementation returns connection coefficients in backward
recurrence order, reverse them before the Horner polynomial evaluation.
The zero sequence has length $0$ and $g=1$.
Reject the trial immediately if $g(0)=0$.

\begin{lemma}[Unconditional lower-rank certificate]\label{lem:lower}
For every $T,u,v$, the finite-sequence complexity in~\eqref{eq:seq}
satisfies $\ell\le\rank T=k$.
\end{lemma}
\begin{proof}
$Tv$ lies in the invariant subspace $\operatorname{im}T$, of
dimension $k$. Cayley--Hamilton on that subspace yields an infinite
recurrence of order at most $k$ and therefore one of that order on
the finite prefix. The minimal prefix length cannot exceed $k$.
This includes singular matrices with nilpotent blocks.
\end{proof}

Set $\widehat r=n-\ell$. Generate at most $n+1$ fresh independent
uniform $b\in K^n$. For each, compute
\[
 y=g(T)b
\]
by Horner's rule. Verify $My=0$ exactly, and retain $y$ only if it
is linearly independent of previously retained verified vectors.
Stop and accept the trial when $\widehat r$ independent verified
vectors have been obtained. If $\widehat r=0$, accept with the empty
basis immediately. If the cap is exhausted, discard the trial and
start another independent one.

Acceptance is sound on \emph{every} random choice. The verified
independent vectors give $\rank M\le n-\widehat r=\ell$;
Lemma~\ref{lem:lower} gives the reverse inequality. Thus acceptance
certifies $\ell=k$ and a complete extension-field kernel basis.
An underestimated rank cannot pass: it asks for more independent
kernel vectors than exist. A bad polynomial action cannot pass
without the explicit kernel and independence checks.

\section{Good-trial probability, probes, and expected work}
Conditional on~\eqref{eq:cyc}--\eqref{eq:nil} being good, random
left and right probes retain the minimum polynomial of the invertible
cyclic $k$-dimensional restriction with failure probability at most
$2k/q$. This follows from the nonzero degree-$k$ controllability and
observability determinant polynomials and Schwartz--Zippel.
The restriction's projected scalar minimum polynomial is determined
exactly by the first $2n$ terms, since its degree is at most $n$;
this is the scalar Wiedemann recurrence fact \cite[Section II]{wiedemann}.
Hence on the combined good event $\mathcal G$,
\[
 \ell=k,\quad g(0)\ne0,\quad
 g(T)=g(0)P_{\ker T\parallel\operatorname{im}T},
\]
and
\begin{equation}\label{eq:bad}
 \Prob(\mathcal G^c)
 \le\frac{n(n-1)/2+r}{q-1}+\frac{2k}{q}
 =O(n^2/q).
\end{equation}
For $k=0$ the statement means $g=1$ and the projection is the identity.

Thus each $g(T)b$ on a good trial is an independent uniform vector
in the $r$-dimensional kernel (up to a nonzero scalar).
The expected number of draws to span it, without a cap, is
\[
 \sum_{j=0}^{r-1}\frac1{1-q^{j-r}}\le r+2,
\]
with zero draws when $r=0$.
After $n+1$ draws, the failure probability is at most
\begin{equation}\label{eq:cap}
 \frac{q^r-1}{q-1}\,q^{-(n+1)}
 \le\frac1{q(q-1)}.
\end{equation}
Indeed a proper span is contained in a hyperplane, and there are
$(q^r-1)/(q-1)$ such hyperplanes, all samples lying in a given one with probability
$q^{-(n+1)}$. Our field size makes each trial succeed with
probability at least $1/2$, uniformly over fixed $M$.

\paragraph{Charging an entire trial.}
The sequence, dot products, and recurrence require
$O(n(w+n))$ field operations, $2n$ applications of $T$, and
$O(n)$ temporary field storage. One polynomial action requires
$O(n(w+n))$ field operations. Kernel verification costs $O(w)$
per candidate. Maintaining independent verified columns costs
$O(nj)$ operations when $j$ columns have already been retained.
On a good trial these costs total in expectation
\[
 O\bigl(n(w+n)(r+1)+nr(r+1)\bigr)
 =O\bigl(n(w+n)(r+1)\bigr).
\]
Even on a bad trial the cap gives the deterministic work bound
$O(n^2(w+n))$. By~\eqref{eq:bad}, its expected contribution is
$O(n^4(w+n)/q)$, absorbed by the displayed good-trial bound.
At all times there are at most $r$ retained vectors, since they
are verified and independent, even on bad trials.

Independent retries satisfy the renewal identity
$\E W=\E W_{\mathrm{trial}}/\Prob(\mathrm{accept})$.
They therefore terminate almost surely with the asserted expected
field work on \emph{every fixed matrix}, including all exceptional
graphs. No rational fallback or high-probability corank-one event
is hidden in this Las Vegas assertion.

\section{Descent to the base field and deterministic verification}
Let $B\in K^{n\times r}$ be an accepted full kernel basis.
Find row indices $J$ with $B_{J,:}$ invertible and compute
\[
 C=B(B_{J,:})^{-1},\qquad C_{J,:}=I_r.
\]
This costs $O(nr^2+r^3)$ field operations, bounded by
$O(n(w+n)(r+1))$, and $O(nr)$ additional storage.

\begin{lemma}[Canonical normalization descends]\label{lem:descent}
Every entry of $C$ belongs to the embedded base field $\F_p$.
\end{lemma}
\begin{proof}
The kernel is defined by $M$ with base-field coefficients, so applying
the coefficientwise Frobenius $x\mapsto x^p$ to $C$ yields another
kernel basis with the same identity on $J$. Restriction to these
pivot coordinates is a bijection from the kernel to $K^r$.
Consequently this normalized basis is unique, and $C^{(p)}=C$.
The fixed elements of this Frobenius in $K$ are precisely $\F_p$.
\end{proof}

In the polynomial representation, check that all positive-degree
coefficients of $C$ vanish, read off the base-field entries, and
verify $MC=0$ in $\F_p$. These are exact checks, not an assumption
that arbitrary extension-field vectors can simply be enumerated as
base-field vectors.

The certificate comprises the irreducible polynomial $h$, $D,u,v$,
$\ell$, the base-field matrix $C$, and its pivot rows $J$.
A verifier checks the representation sizes and bounds on $q$, tests
irreducibility, recomputes~\eqref{eq:seq} and its exact BM length,
and checks
\[
 r=n-\ell,\qquad MC=0,\qquad C_{J,:}=I_r.
\]
These give the lower and upper rank bounds deterministically, without
trusting the random preconditioner or the polynomial projection.
Verification needs $O(n(w+n)+wr+r^2)$ field/base operations, within
the theorem's bound. It does not recompute a dense elimination.

\section{Field construction, bit work, and storage}
The extension degree is $s=O_p(\log(n+1))$. Draw monic degree-$s$
polynomials over $\F_p$ and test irreducibility deterministically.
The standard irreducible-polynomial count gives success probability
at least $1/(2s)$, so the expected number of draws is $O(s)$.
One may use Rabin's exact criterion \cite{rabin}, or the elementary
version checking
\[
 \gcd\bigl(h,X^{p^i}-X\bmod h\bigr)=1
 \quad\text{for all }1\le i\le\lfloor s/2\rfloor.
\]
A reducible degree-$s$ polynomial has an irreducible factor of degree
at most $s/2$, so this version is sufficient and necessary.
Successive modular Frobenius powers and classical polynomial gcd
give polynomial work in $s$, hence polylogarithmic bit work in $n$.
Every accepted field construction is correct; randomness only
controls its running time.

Represent elements by $s$ base-field coefficients. Addition,
multiplication modulo $h$, and inversion by extended Euclid each
cost a polynomial in $s$, hence $\widetilde O_p(1)$ bit operations.
Independent uniform field choices use $O(\log q)$ random bits in
expectation per element; rejecting zero for diagonal entries has
bounded expected overhead. This charges all field construction,
randomness, polynomial actions, and basis arithmetic in
Theorem~\ref{thm:kernel}. Workspace is $O(n(r+1)+w)$ field/base
elements; each has $O_p(\log(n+1))$ bits. Retry-counter storage can
be treated separately and has finite expected size.

\section{Exact partition listing from a single corank moment}
Let $M_{\mathbb Z}=A-\lambda I$ for a simple graph of maximum degree
$d$ and a known integer $|\lambda|\le d$, and let $M$ be its reduction modulo
the fixed odd prime $p$. Compute the full base-field kernel by
Theorem~\ref{thm:kernel}. Enumerate all $p^r$ coefficient vectors
$a\in\F_p^r$ and their images $Ca$.
Retain only those with last coordinate $1$ and every coordinate in
$\{1,-1\}$. Lift these residues uniquely to integer signs and
check exactly
\[
 \one^T\sigma=0,\qquad M_{\mathbb Z}\sigma=0.
\]
Return every passing sign vector, with the orientation $\sigma_n=1$.
The list is complete on every input: every true balanced integer
sign kernel vector reduces to a vector of this finite-field kernel.
Its integer verification eliminates all modular artifacts.

\paragraph{The enumeration factor is charged without an extra $r$.}
Use a base-$p$ odometer or Gray ordering and maintain the current
image $Ca$. In the full run, the number of digit changes is at most
$p^r\sum_{j\ge0}p^{-j}=O_p(p^r)$.
Each digit change adds a known basis column in $O(n)$ base operations.
All sign, balance, and integer degree checks cost $O_d(n)$ per
candidate. The full enumeration therefore costs
\[
 O_{d,p}(n p^r)
\]
bit operations, including the list entries. Recomputing every image
by a fresh $nr$ multiplication would be an invalid expectation
shortcut: $\E p^r<\infty$ need not bound $\E(rp^r)$.

\begin{corollary}[Moment-to-full-list theorem]\label{cor:moment}
Suppose a bounded-degree random-graph model satisfies, for one fixed
odd prime $p$,
\[
 \sup_n\E\bigl[p^{\dim\ker_{\F_p}(A-\lambda I)}\bigr]\le C<\infty.
\]
Then the preceding algorithm lists all balanced integer sign kernels
of $A-\lambda I$, up to global exchange, in unconditional expected
$\widetilde O_{d,p,C}(n^2)$ bit work. It has no algorithmic error
and does not assume real or modular corank one. Its expected working
storage is $\widetilde O_{d,p,C}(n)$ bits and its expected output
size is $O_{d,p,C}(n)$ bits.
\end{corollary}
\begin{proof}
The per-input conditional expected work is
\[
 \widetilde O_{d,p}\bigl(n^2(r+1)+n p^r\bigr).
\]
Since $r+1\le O_p(p^r)$, the stated raw moment bounds its expectation
by $\widetilde O(n^2)$ and gives the storage/output bounds as well.
This averages a valid all-input Las Vegas bound, so exceptional
graphs incur no unjustified dense fallback cost.
\end{proof}

\paragraph{Exactly what is assumed for the RSBM application.}
The root research packet proposes a bounded raw kernel moment for
$\lambda=a-b$, a fixed prime $p>b$, internal degree $a\ge3$ on
all admissible block sizes, and separately the cubic case $(a,b)=(2,1)$.
Corollary~\ref{cor:moment} uses that moment as an input.
If it is established directly in the configuration model and
simplicity has probability bounded below, conditioning gives
\[
 \E[p^r\mid\mathrm{simple}]
 \le\frac{\E p^r}{\Prob(\mathrm{simple})}=O(1).
\]
Qualitative contiguity alone would not justify this raw-moment transfer.
The present note audits the algorithmic implication and does not
independently certify every all-target configuration-mixing proof.

For a regular graph, a sign vector satisfying
$A\sigma=(a-b)\sigma$ has exactly $a$ same-side and $b$ opposite-side
neighbors at every vertex, because its total degree is $a+b$.
Thus the verified list is exactly the list of admissible balanced
partitions for those parameters. On a planted input it is nonempty.
Returning the sole partition when the list has size one, and the
complete list when larger, leaves ambiguity as a statistical property
of the input, rather than a computational failure on extra eigenspaces.
No ambiguity probability or priority claim is supplied by this note.

\paragraph{Validation and remaining dependency.}
The accompanying exact prototype checks zero and full rank, several
coranks, a symmetric nilpotent input in characteristic three, exact
extension construction, base-field descent, and altered certificates.
It recovers the unique admissible partition of the cubic $N=8$
fixture with internal cycles $(0,1,2,3,4,5,6,7)$ and
$(5,6,7,0,4,1,2,3)$ and identity matching, despite that fixture's
real planted eigenspace having dimension two.
It also lists both admissible partitions in a separate ambiguous
$12$-vertex example, agreeing with exhaustive sign search.

The model-level expected-time conclusion still requires the stated
raw kernel moment for the \emph{actual simple-graph distribution}.
A high-probability simple-real-kernel statement alone is insufficient.
The result here uses a known integer $\lambda$; unconditional expected
cost for enumerating other integer shifts would require the relevant
moments for those shifts or a separately charged filtering argument.

\begin{thebibliography}{9}
\bibitem{chen} L.~Chen, W.~Eberly, E.~Kaltofen, B.~D.~Saunders,
W.~J.~Turner, G.~Villard,
\emph{Efficient matrix preconditioners for black box linear algebra},
Linear Algebra and its Applications 343--344 (2002), 119--146.
Theorems 4.2 and 4.5.
\url{https://kaltofen.math.ncsu.edu/bibliography/02/CEKSTV02.pdf}.
\bibitem{wiedemann} D.~H.~Wiedemann,
\emph{Solving Sparse Linear Equations Over Finite Fields},
IEEE Transactions on Information Theory 32 (1986), 54--62.
\url{https://doi.org/10.1109/TIT.1986.1057137}.
\bibitem{rabin} M.~O.~Rabin,
\emph{Probabilistic Algorithms in Finite Fields},
SIAM Journal on Computing 9 (1980), 273--280.
\url{https://doi.org/10.1137/0209024}.
\end{thebibliography}
\end{document}
